% hole_nyquist.tex -- round 410, one lemma in two halves:
% the hole-Nyquist defect after r407/r408/r409.
% N: TESTPOLY_REFUTED / DMIN_CENSUS
% S: FOURIER_REFUTED / BERNSTEIN_REFUTED / SEQ_REDUCED
% Builds with pdflatex.
% Companion: verify_hole_nyquist.py
% Discovery: experiments/tfpt-discovery/hole_nyquist_probe.py
% NO RH CLAIM.  Research documentation, not a theorem of RH.
\documentclass[11pt]{article}

\usepackage[a4paper,margin=2.7cm]{geometry}
\usepackage{amsmath,amssymb,amsthm}
\usepackage{microtype}
\usepackage{booktabs}
\usepackage{xcolor}
\usepackage[colorlinks=true,linkcolor=blue!50!black,urlcolor=blue!50!black]{hyperref}

\newtheorem{lemma}{Lemma}
\newtheorem{prop}{Proposition}
\theoremstyle{remark}
\newtheorem{remark}{Remark}

\title{\bfseries The hole-Nyquist defect\\[2mm]
\large Test polynomials for $d_i\ge 1$, and why the
off-diagonal mass is not a Fourier or Bernstein theorem}
\author{Stefan Hamann}
\date{August 29, 2026}

\begin{document}
\maketitle

\begin{abstract}\noindent
After r407 ($P1\iff\lambda_2(C)\ge 1$) and r408
($d_{\min}\ge 1$ is the necessary half, the exception mode
is hole-Nyquist), this note attacks both halves.
Half~N: the reproducing-kernel test-polynomial bound with
the $Y$-Lagrange interpolant (degree $|Y|-1<n$) does
\emph{not} prove $d_i\ge 1$ ($b_{\max}=0.071$, zero holes
cleared).  Half~S: $C$ is not banded in the hole-DFT
(off-fraction $0.968$); $\|T_0 v_{\mathrm{Nyq}}\|<\|T_0\|$;
the defect is born at $k=37$ along the $\theta$-prefix and
stays unique.  No RH claim.
\end{abstract}

\medskip\noindent
\textbf{Status.}\enspace
Lemmas~\ref{lem:rk}--\ref{lem:rayleigh} are proved over
$\mathbb{Q}$.
The test-polynomial SATZ-candidate, the Fourier-block
candidate, and the Bernstein candidate are
\textbf{refuted}.
$d_i\ge 1$ on the source family is a census.
The sequential birth is a named reduction, not a theorem.
No $L^*$ claim.  No $R^\dagger$ claim.  No RH claim.

\section{The object}

$C=BB^T$ is the $\mu$-only dual CD Gram on $Y$ (r407),
\[
  d_i \;=\; C_{ii}
  \;=\; u^\vee(y_i)\,K_n^\mu(y_i,y_i).
\]
Rayleigh gives $d_{\min}<1\Rightarrow\lambda_{\min}<1$.
r408: source $d_{\min}\in[1.097,1.656]$; scramble $0.37$;
permute $0.15$ (same $Y$).

\section{Half N: the test polynomial}

\begin{lemma}[Reproducing kernel]\label{lem:rk}
For $p\in\mathcal{P}_n$,
$K_n^\mu(y,y)\ge |p(y)|^2/\|p\|_{L^2(\mu)}^2$.
On three unit atoms, $n=1$, $K=1/h_0=1/3$.
\end{lemma}

\begin{lemma}[Degree obstruction]\label{lem:deg}
If $\deg p\ge n$ the inequality may fail:
$p(t)=t^2$ on $\{-1,0,1\}$, $n=1$, gives
$|p(1)|^2/\|p\|^2=1/2>1/3=K$.
\end{lemma}

\begin{lemma}[$C_{ii}=u^\vee K$]\label{lem:cii}
$C=BB^T$ with $B_{ki}=\sqrt{u^\vee(y_k)}\,P_i(y_k)$ implies
$C_{ii}=u^\vee(y_i)\sum_{j<n}P_j(y_i)^2$.
\end{lemma}

\begin{lemma}[Rayleigh]\label{lem:rayleigh}
$\lambda_{\min}(C)\le\min_i C_{ii}$.
\end{lemma}

The candidate $p_{y_i}$ is the $Y$-Lagrange basis
$\ell_i$ (zeros on $Y\setminus\{y_i\}$, $\ell_i(y_i)=1$).
On FRAME-A, $\deg\ell_i=103<n=181$, so
Lemma~\ref{lem:rk} applies, and
\[
  d_i \;\ge\; b_i
  \;:=\; u^\vee(y_i)\big/\textstyle\sum_{x\in X}u^\vee(x)\,\ell_i(x)^2.
\]

\begin{prop}[Lagrange does not close]\label{prop:lag}
On $kz=9$, $b_{\max}=0.071$, $\#\{b_i\ge 1\}=0/104$,
$\min(d_i/b_i)=94$ (no RK violation).
Local Lagrange is worse.
The Chebyshev-CD test of degree $n-1$ clears $81/104$
holes and leaves $b_{\min}=4.9\cdot 10^{-4}$.
\end{prop}

So Half~N is \textbf{not} a SATZ by this family of explicit
$p$.  $d_i\ge 1$ remains a source census (core-$42$,
$d_{\min}\in[1.097,1.656]$).

Permute (same $Y$) drops $d_{\min}$ to $0.148$ by a
\emph{Christoffel} collapse: $K_n^\mu(y,y)=0.148$ at a hole
with $u^\vee=1$.  The test-poly bounds stay tiny on
scramble and permute --- they do not exhibit the break.
Live and dead $\chi_3$ keep $d_{\min}>1$.

\section{Half S: one Nyquist mode}

\begin{prop}[Fourier is not a block]\label{prop:ft}
The hole-DFT of $C$ has off-diagonal energy fraction
$0.968$; an $8$-diagonal band holds only $0.18$ of the
energy.  The Nyquist bin is not a spectral atom.
\end{prop}

\begin{prop}[Bernstein is not tight]\label{prop:bern}
$\|T_0 v_{\mathrm{Nyq}}\|=0.656<\|T_0\|=1.080$
(r409 graph operator).
The defect eigenvector has hole-Nyquist overlap $0.386$,
not $1$.
A Markov estimate of the required slope over the minimal
$x$-gap is $0.14$ of the $n^2$ bound --- not critical.
\end{prop}

\begin{prop}[Sequential birth]\label{prop:seq}
Rebuilding $C$ along the $\theta$-prefix of $Y$:
$k=32$ has $n_C=0$, $C_{\min}=1.00014$;
$k=37$ has $n_C=1$, $C_{\min}=0.99991$;
$k=104$ has $n_C=1$, $\lambda_2=1.00018$.
The defect is born once and does not split.
The flip index depends on order (a random prefix flips
later).  This reduces Half~S to a one-path census, not a
Cauchy theorem (rebuilt $C$ is not a principal submatrix).
\end{prop}

\begin{remark}
P1 remains equivalent to $\lambda_2(C)\ge 1$.
Neither half proves the right-hand side.
Unconditional SATZ is Lemmas~\ref{lem:rk}--\ref{lem:rayleigh}.
$1010$ at $S=20$ has $\deg\ell>|Y|$ vs.\ $n$ and violates
RK ($d<b$), as predicted by Lemma~\ref{lem:deg}.
\end{remark}

\section*{Kill table (machine)}

\begin{center}
\begin{tabular}{@{}lcc@{}}
\toprule
mutant & effect & what it kills \\
\midrule
$Y$-Lagrange as SATZ & $b_{\max}=0.071$ & Half N proof \\
Chebyshev-CD $m=n$ & $81/104$ & remaining holes \\
scramble & $d_{\min}=0.37$ & Half N on that world \\
$\Lambda$-permute & $K\to 0.148$ & same $Y$, weights \\
$1010$ $S=20$ & $\deg>n$, $d<b$ & RK hypothesis \\
flip $+10$ & $n_C=8$ & density \\
depth $n+1$ & $C_{\min}\to 1$ & razor \\
\bottomrule
\end{tabular}
\end{center}

\end{document}
